\documentclass[11pt,a4paper]{article}
\usepackage[margin=1in]{geometry}
\usepackage{amsmath,amssymb}
\usepackage{enumitem}
\usepackage{array}
\usepackage{longtable}

\setlength{\parindent}{0pt}
\setlength{\parskip}{6pt}

\title{\textbf{CAT 2025 Quant -- CAT-Hardened Practice Paper}\\[4pt]
\large Hard-End Simulation (15 Questions)}
\author{}
\date{}

\begin{document}
\maketitle

\section*{Instructions}
\begin{itemize}[leftmargin=*]
  \item This paper contains 15 questions of the Quantitative Ability section.
  \item Total questions: 15. Multiple-Choice (MCQ): 9. Type-In-The-Answer (TITA): 6.
  \item Suggested time: 55 minutes.
  \item MCQ: each has exactly one correct option. TITA: enter the exact numerical value.
  \item No calculators permitted. Rough work may be done on a separate sheet.
  \item Topic, sub-topic, difficulty, and type are shown next to each question for self-diagnosis.
\end{itemize}

\hrule

\section*{Question Paper}

\textbf{Q1.} \textit{[Arithmetic -- Profit \& Loss, Medium-Hard, MCQ]}\\
A retailer marks every item in his store at 60\% above its cost price. As part of a promotion, on every third item that is sold he gives a discount of $d\%$ on the marked price, while the remaining items are sold at the marked price. Over a block of 9 items sold under this scheme, his overall profit is exactly 40\%. The value of $d$ is:
\begin{enumerate}[label=(\Alph*)]
  \item 30
  \item 32.5
  \item 35
  \item 37.5
\end{enumerate}

\textbf{Q2.} \textit{[Arithmetic -- Time \& Work / Pipes, Hard, TITA]}\\
Pipes $A$ and $B$ can fill an empty tank in 10 hours and 15 hours respectively, while pipe $C$ can empty a full tank in 30 hours. Starting at time $0$ with an empty tank, the pipes are operated as follows: in the 1st hour only $A$ is open; in the 2nd hour $A$ and $B$ are open; in the 3rd hour $A$, $B$ and $C$ are all open. This 3-hour pattern then repeats. The total time, in hours, after which the tank first becomes full is \rule{2cm}{0.4pt}.

\textbf{Q3.} \textit{[Arithmetic -- Time, Speed \& Distance, Hard, MCQ]}\\
Two cars $X$ and $Y$ travel between cities $P$ and $Q$ along the same straight road of length $D$ km. Car $X$ leaves $P$ at 9:00~a.m. heading towards $Q$ at 50 km/h. Car $Y$ leaves $Q$ at 9:45~a.m. heading towards $P$ at 40 km/h. On reaching the other city, each car immediately turns around and continues at the same speed. They meet for the second time at exactly 12:00~noon. The value of $D$ (in km) is:
\begin{enumerate}[label=(\Alph*)]
  \item 75
  \item 80
  \item 90
  \item 100
\end{enumerate}

\textbf{Q4.} \textit{[Arithmetic -- Mixtures \& Alligation, Medium-Hard, MCQ]}\\
Two solutions $A$ and $B$ contain alcohol and water in the ratios $5:3$ and $7:5$ respectively. $A$ and $B$ are mixed in some ratio to produce a solution $C$ in which alcohol and water are in the ratio $8:5$. The solution $C$ is then mixed with pure water in the ratio $13:k$ (by volume) so that the final solution has alcohol concentration exactly 40\%. The value of $k$ is:
\begin{enumerate}[label=(\Alph*)]
  \item 5
  \item 6
  \item 7
  \item 8
\end{enumerate}

\textbf{Q5.} \textit{[Arithmetic -- Time \& Work, Hard, TITA]}\\
Working alone, $A$ and $B$ can complete a certain piece of work in 20 days and 30 days respectively. They are assigned to work together every day, except that $A$ takes leave on every 5th day (i.e.\ on days $5, 10, 15, \dots$); on such days only $B$ works. Starting on Day 1, the smallest positive integer $N$ such that the work is completely finished by the end of day $N$ is \rule{2cm}{0.4pt}.

\textbf{Q6.} \textit{[Algebra -- Quadratics / Integer feasibility, Hard, MCQ]}\\
The number of integer values of $k$ for which the equation
\[
x^{2} - (k+3)\,x + (4k+1) = 0
\]
has both roots as positive integers is:
\begin{enumerate}[label=(\Alph*)]
  \item 0
  \item 1
  \item 2
  \item 3
\end{enumerate}

\textbf{Q7.} \textit{[Algebra -- Inequalities with parameter, Hard, MCQ]}\\
The number of integer values of $a$ for which
\[
(a-2)\,x^{2} + (a-2)\,x + 1 > 0
\]
holds for every real number $x$ is:
\begin{enumerate}[label=(\Alph*)]
  \item 2
  \item 3
  \item 4
  \item 5
\end{enumerate}

\textbf{Q8.} \textit{[Algebra -- Functional equations, Hard, TITA]}\\
A function $f:\mathbb{R} \to \mathbb{R}$ satisfies
\[
f(x+y) + f(x-y) = 2\,f(x)\,f(y) \quad \text{for all real } x, y,
\]
with $f(1) = \tfrac{1}{2}$. The value of $f(2025) + f(2026)$ is \rule{2cm}{0.4pt}.

\textbf{Q9.} \textit{[Algebra -- Symmetric polynomials / Optimisation, Very Hard, MCQ]}\\
Let $a, b, c$ be non-negative real numbers satisfying $a + b + c = 6$ and $ab + bc + ca = 9$. The maximum possible value of $abc$ is:
\begin{enumerate}[label=(\Alph*)]
  \item 2
  \item 3
  \item 4
  \item 6
\end{enumerate}

\textbf{Q10.} \textit{[Number System -- Divisibility, Hard, TITA]}\\
The number of positive integers $n$ with $1 \le n \le 1000$ such that $(n+7)$ divides $(n^{3} + 1)$ is \rule{2cm}{0.4pt}.

\textbf{Q11.} \textit{[Number System -- Remainders / CRT, Medium-Hard, MCQ]}\\
Let $N$ be the smallest positive integer such that when $N$ is divided by 7, 11, and 13, the remainders are 5, 9 and 11, respectively. The sum of the digits of $N$ equals:
\begin{enumerate}[label=(\Alph*)]
  \item 18
  \item 21
  \item 24
  \item 27
\end{enumerate}

\textbf{Q12.} \textit{[Number System -- Base representation / Factorials, Hard, MCQ]}\\
Let $N_{10}$ be the number of trailing zeros when $2025!$ is written in base 10, and let $N_{12}$ be the number of trailing zeros when $2025!$ is written in base 12. The value of $N_{10} - N_{12}$ is:
\begin{enumerate}[label=(\Alph*)]
  \item $-503$
  \item $-500$
  \item $500$
  \item $503$
\end{enumerate}

\textbf{Q13.} \textit{[Geometry -- Triangles / Cevians, Hard, MCQ]}\\
In triangle $ABC$, point $D$ lies on side $BC$ with $BD : DC = 2 : 3$, and point $E$ lies on side $AC$ with $AE : EC = 1 : 2$. Segments $AD$ and $BE$ intersect at $P$. The ratio of the area of triangle $APE$ to the area of triangle $ABC$ is:
\begin{enumerate}[label=(\Alph*)]
  \item $1 : 6$
  \item $1 : 9$
  \item $1 : 12$
  \item $2 : 15$
\end{enumerate}

\textbf{Q14.} \textit{[Geometry -- Triangle / Incircle, Very Hard, TITA]}\\
In triangle $ABC$, $AB = 14$, $BC = 13$ and $CA = 15$. Let $M$ be the midpoint of side $BC$, and let $P$ be the point at which the incircle of triangle $ABC$ touches side $AB$. The value of $4 \cdot MP^{2}$ is \rule{2cm}{0.4pt}.

\textbf{Q15.} \textit{[Mensuration -- Cone \& inscribed spheres, Very Hard, TITA]}\\
A right circular cone has its apex at the top and opens downward, resting on its circular base. The half-vertex angle $\alpha$ of the cone (i.e.\ the angle between its axis and any slant edge) satisfies $\sin \alpha = \tfrac{3}{5}$. A sphere $S_{1}$ of radius 3 cm is placed inside the cone, tangent to its lateral surface and to its horizontal base. A second sphere $S_{2}$ is placed inside the cone, above $S_{1}$, tangent to the lateral surface of the cone and externally tangent to $S_{1}$. The radius of $S_{2}$ (in cm) is \rule{2cm}{0.4pt}.

\vspace{8pt}\hrule

\section*{Answer Key}

\begin{center}
\begin{tabular}{|c|c|c|c|}
\hline
\textbf{Q.No.} & \textbf{Type} & \textbf{Answer} & \textbf{Difficulty} \\
\hline
1  & MCQ  & (D) 37.5       & Medium-Hard \\
2  & TITA & 7.6            & Hard \\
3  & MCQ  & (B) 80         & Hard \\
4  & MCQ  & (C) 7          & Medium-Hard \\
5  & TITA & 14             & Hard \\
6  & MCQ  & (B) 1          & Hard \\
7  & MCQ  & (C) 4          & Hard \\
8  & TITA & $-1.5$         & Hard \\
9  & MCQ  & (C) 4          & Very Hard \\
10 & TITA & 8              & Hard \\
11 & MCQ  & (D) 27         & Medium-Hard \\
12 & MCQ  & (A) $-503$     & Hard \\
13 & MCQ  & (B) $1:9$      & Hard \\
14 & TITA & 193            & Very Hard \\
15 & TITA & 0.75           & Very Hard \\
\hline
\end{tabular}
\end{center}

\hrule
\section*{Detailed Solutions}

\textbf{Solution 1.} \textit{[Profit on a block of items]}\\
Let cost price of each item be $100$ units. Marked price $= 160$.
\begin{itemize}[leftmargin=*]
  \item Out of every 9 items: 6 are sold at MP $= 160$ and 3 are sold at $160(1 - d/100)$.
  \item Total cost over 9 items $= 900$.
  \item Total revenue $= 6(160) + 3 \cdot 160\!\left(1 - \tfrac{d}{100}\right) = 960 + 480\!\left(1 - \tfrac{d}{100}\right)$.
\end{itemize}
Overall profit $= 40\% \Rightarrow$ revenue $= 1260$.
\[
960 + 480 - 4.8\,d = 1260 \;\Longrightarrow\; 4.8\,d = 180 \;\Longrightarrow\; d = 37.5.
\]
\textbf{Answer: (D) 37.5.}

\textit{Why CAT-level:} Tempting first move is to apply a simple discount-on-MP formula on a single item; the trap is the weighted base of 9 items.

\vspace{4pt}
\textbf{Solution 2.} \textit{[Cyclic operation of pipes]}\\
Hour-wise fill rates:
\[
\text{Hr 1:} \tfrac{1}{10},\quad \text{Hr 2:} \tfrac{1}{10}+\tfrac{1}{15}=\tfrac{1}{6},\quad \text{Hr 3:} \tfrac{1}{10}+\tfrac{1}{15}-\tfrac{1}{30}=\tfrac{2}{15}.
\]
Per 3-hour cycle: $\tfrac{1}{10}+\tfrac{1}{6}+\tfrac{2}{15} = \tfrac{3+5+4}{30} = \tfrac{12}{30} = \tfrac{2}{5}$.
After 2 cycles (6 hours): $\tfrac{4}{5}$ filled, $\tfrac{1}{5}$ remaining.

\textit{Hour 7} (only $A$): fills $\tfrac{1}{10}$. Total $= \tfrac{4}{5}+\tfrac{1}{10}=\tfrac{9}{10}$; remaining $= \tfrac{1}{10}$.

\textit{Hour 8} ($A$ and $B$): fill-rate $\tfrac{1}{6}$ per hour $> \tfrac{1}{10}$. Time to finish in Hour 8 $= \dfrac{1/10}{1/6} = \tfrac{3}{5}$ hour.

Total time $= 7 + \tfrac{3}{5} = \tfrac{38}{5} = 7.6$ hours.
\textbf{Answer: 7.6.}

\textit{Why CAT-level:} The combined rate per cycle is positive, but the third hour has a negative contribution. A naive LCM/cycle-only count misses the threshold-crossing inside the partial cycle.

\vspace{4pt}
\textbf{Solution 3.} \textit{[Bouncing cars, asymmetric start]}\\
Take 9:00~a.m.\ as time $t=0$. Then $Y$ starts at $t=0.75$.
\begin{itemize}[leftmargin=*]
  \item $X$ reaches $Q$ at $t = D/50$, then turns back.
  \item $Y$ reaches $P$ at $t = 0.75 + D/40$, then turns back.
\end{itemize}
For the second meeting to occur with both cars having bounced once each:
\[
\text{Position of } X = 2D - 50t,\qquad \text{Position of } Y = 40(t - 0.75) - D = 40t - 30 - D.
\]
Setting them equal:
\[
2D - 50t = 40t - 30 - D \;\Longrightarrow\; 3D + 30 = 90t \;\Longrightarrow\; D = 30t - 10.
\]
At $t=3$: $D = 80$.
Verification at $t=3$: $X$ is at $80 - 50(3-1.6) = 10$; $Y$ is at $40(3-2.75) = 10$. They coincide at $10$ km from $P$.
\textbf{Answer: (B) 80.}

\textit{Why CAT-level:} The standard ``combined distance $= 3D$ at second meeting'' shortcut breaks because $Y$ starts later. Method selection (bouncing analysis vs.\ symmetric shortcut) is the trap.

\vspace{4pt}
\textbf{Solution 4.} \textit{[Two-stage mixing]}\\
Alcohol fractions: $A = 5/8$, $B = 7/12$. Target for $C$: $8/13$.

\textbf{Stage 1.} If $A : B = p : q$, then
\[
\frac{\tfrac{5p}{8} + \tfrac{7q}{12}}{p+q} = \tfrac{8}{13}
\Longrightarrow 13(15p + 14q) = 192 \cdot 24 \cdot \tfrac{p+q}{24} \cdot \tfrac{24}{1}.
\]
Working it cleanly: $13(15p + 14q) = 192(p+q) \Rightarrow 3p = 10q$, i.e.\ $p:q = 10:3$. (This step is only to confirm $C$ is achievable; the alcohol fraction of $C$ is fixed at $8/13$.)

\textbf{Stage 2.} Take 13 units of $C$: alcohol $= 8$, water $= 5$. Add $k$ units of water. Final alcohol fraction $= 40\% = 2/5$:
\[
\frac{8}{13 + k} = \frac{2}{5} \;\Longrightarrow\; 13 + k = 20 \;\Longrightarrow\; k = 7.
\]
\textbf{Answer: (C) 7.}

\textit{Why CAT-level:} Two linked mixings. Starting with alligation on $C$ + water is the elegant path; mixing $A$ + $B$ explicitly is the false-start consumer of time.

\vspace{4pt}
\textbf{Solution 5.} \textit{[Work with leave schedule]}\\
Daily rates: $A = \tfrac{1}{20},\ B = \tfrac{1}{30}$. Joint rate (no leave) $= \tfrac{1}{12}$. On leave days only $B$: rate $\tfrac{1}{30}$.

Cumulative work at end of each day (in units of $\tfrac{1}{60}$):
\[
\begin{array}{c|cccccccccccccc}
\text{Day} & 1 & 2 & 3 & 4 & 5 & 6 & 7 & 8 & 9 & 10 & 11 & 12 & 13 & 14\\
\hline
\text{Cum.} & 5 & 10 & 15 & 20 & 22 & 27 & 32 & 37 & 42 & 44 & 49 & 54 & 59 & 64
\end{array}
\]
(Each non-leave day adds $5$; leave days $5, 10$ add $2$.) The total work corresponds to $60$. The cumulative crosses $60$ during Day 14. Hence the work is completely finished by the end of Day 14.
\textbf{Answer: $N = 14$.}

\textit{Why CAT-level:} The trap is to compute average daily rate over a 5-day block, miss the residual, and undercount.

\vspace{4pt}
\textbf{Solution 6.} \textit{[Integer roots]}\\
Let the positive integer roots be $p, q$.
By Vieta, $p+q = k+3$ and $pq = 4k+1$. Substituting $k = p+q-3$:
\[
pq = 4(p+q-3) + 1 = 4p + 4q - 11.
\]
Rearrange to apply factor-shifting:
\[
pq - 4p - 4q + 16 = 5 \;\Longrightarrow\; (p-4)(q-4) = 5.
\]
With $p, q \in \mathbb{Z}^{+}$, the factor pairs giving both $p, q \ge 1$ are $(p-4, q-4) \in \{(1,5),(5,1)\}$, yielding $(p,q) \in \{(5,9),(9,5)\}$. Both correspond to $k = 5+9-3 = 11$.
Hence exactly \emph{one} integer value of $k$ works.
\textbf{Answer: (B) 1.}

\textit{Why CAT-level:} Direct discriminant-perfect-square attempt is messy; the elegant route requires Simon-style factor shifting.

\vspace{4pt}
\textbf{Solution 7.} \textit{[Inequality for all real $x$]}\\
\textbf{Case 1:} $a - 2 = 0$, i.e.\ $a = 2$. The inequality becomes $1 > 0$ -- always true.

\textbf{Case 2:} $a - 2 \neq 0$. For the quadratic to be positive for all real $x$ we need $a-2 > 0$ \emph{and} discriminant $< 0$:
\[
(a-2)^{2} - 4(a-2) < 0 \;\Longrightarrow\; (a-2)(a-6) < 0 \;\Longrightarrow\; 2 < a < 6.
\]
Combining with $a-2 > 0$ gives $2 < a < 6$.

Total set of valid $a$: $\{2\} \cup (2, 6) = [2, 6)$. Integer values: $\{2, 3, 4, 5\}$.
\textbf{Answer: (C) 4.}

\textit{Why CAT-level:} The natural reflex (set discriminant $< 0$) silently kills the $a = 2$ degenerate case.

\vspace{4pt}
\textbf{Solution 8.} \textit{[D'Alembert-type functional equation]}\\
Put $x = y = 0$: $2f(0) = 2f(0)^{2}$, so $f(0) \in \{0, 1\}$. Trial with $f(0) = 0$ leads to a contradiction (computing $f(4)$ via two routes), so $f(0) = 1$.

Now put $y = 1$ to obtain a recurrence:
\[
f(x+1) + f(x-1) = 2f(x)\,f(1) = f(x) \;\Longrightarrow\; f(x+1) = f(x) - f(x-1).
\]
With $f(0) = 1,\ f(1) = \tfrac{1}{2}$:
\[
f(0), f(1), \dots, f(5) \;=\; 1,\ \tfrac{1}{2},\ -\tfrac{1}{2},\ -1,\ -\tfrac{1}{2},\ \tfrac{1}{2},
\]
and $f(6) = 1 = f(0)$ -- periodic with period 6.

Since $2025 = 6 \cdot 337 + 3$ and $2026 \equiv 4 \pmod{6}$:
\[
f(2025) = f(3) = -1,\qquad f(2026) = f(4) = -\tfrac{1}{2}.
\]
\[
f(2025) + f(2026) = -\tfrac{3}{2}.
\]
\textbf{Answer: $-1.5$.}

\textit{Why CAT-level:} Brute trigonometric guess ($f(x) = \cos(\pi x/3)$) works but takes much longer than the integer-only recurrence.

\vspace{4pt}
\textbf{Solution 9.} \textit{[Max of $abc$ with fixed $e_1, e_2$]}\\
The numbers $a, b, c$ are the roots of the cubic
\[
t^{3} - 6t^{2} + 9t - p = 0, \quad \text{where } p = abc.
\]
Define $g(t) = t(t-3)^{2} = t^{3} - 6t^{2} + 9t$. The cubic becomes $g(t) = p$. We need three real non-negative roots.

$g'(t) = 3(t-1)(t-3)$: local max $g(1) = 4$, local min $g(3) = 0$. For $0 \le p \le 4$, $g(t) = p$ has three real roots. For $0 < p < 4$, all three roots lie in $(0, \infty)$ (one in $(0,1)$, one in $(1,3)$, one in $(3, \infty)$). At the boundary $p = 4$, the roots are $1, 1, 4$, which satisfy $1+1+4 = 6$ and $1\!\cdot\!1 + 1\!\cdot\!4 + 1\!\cdot\!4 = 9$. For $p > 4$, only one root is real.

Hence the maximum of $abc$ is $4$, attained at $(a, b, c) = (1, 1, 4)$ (and permutations).
\textbf{Answer: (C) 4.}

\textit{Why CAT-level:} The Newton-Vieta linkage between $abc$ and a cubic in one variable is non-obvious; AM-GM gives only weak bounds here.

\vspace{4pt}
\textbf{Solution 10.} \textit{[Divisibility shift]}\\
By the remainder theorem, $n^{3} + 1$ leaves remainder $(-7)^{3} + 1 = -342$ when divided by $n + 7$. Hence
\[
(n+7) \mid (n^{3} + 1) \iff (n+7) \mid 342.
\]
Now $342 = 2 \cdot 3^{2} \cdot 19$ has $12$ positive divisors:
$1, 2, 3, 6, 9, 18, 19, 38, 57, 114, 171, 342$.
We need $n \ge 1$ and $n \le 1000$, i.e.\ $8 \le n+7 \le 1007$. Divisors of $342$ in this range: $9, 18, 19, 38, 57, 114, 171, 342$ -- a total of $8$ values.
\textbf{Answer: 8.}

\textit{Why CAT-level:} The transformation $n+7 \mid n^{3}+1 \iff n+7 \mid 342$ is the non-obvious step; brute enumeration over $n$ is hopeless.

\vspace{4pt}
\textbf{Solution 11.} \textit{[Hidden congruence pattern]}\\
Observe:
\[
N \equiv 5 \pmod{7} \iff N \equiv -2 \pmod{7},\quad
N \equiv 9 \pmod{11} \iff N \equiv -2 \pmod{11},\quad
N \equiv 11 \pmod{13} \iff N \equiv -2 \pmod{13}.
\]
Hence $N \equiv -2 \pmod{\mathrm{lcm}(7, 11, 13)} = -2 \pmod{1001}$, so $N = 1001k - 2$ for some $k \ge 1$. The smallest positive value is $N = 999$.
Sum of digits $= 9 + 9 + 9 = 27$.
\textbf{Answer: (D) 27.}

\textit{Why CAT-level:} A direct CRT setup is correct but slow. The shortcut requires noticing each remainder is ``two less than the modulus.''

\vspace{4pt}
\textbf{Solution 12.} \textit{[Trailing zeros in two bases]}\\
\textbf{Base 10.} Number of trailing zeros equals the power of $5$ in $2025!$ (since the power of $2$ is far larger):
\[
N_{10} = \left\lfloor \tfrac{2025}{5} \right\rfloor + \left\lfloor \tfrac{2025}{25} \right\rfloor + \left\lfloor \tfrac{2025}{125} \right\rfloor + \left\lfloor \tfrac{2025}{625} \right\rfloor = 405 + 81 + 16 + 3 = 505.
\]

\textbf{Base 12.} Since $12 = 2^{2}\!\cdot\!3$, trailing zeros equal the largest $k$ such that $2^{2k} \cdot 3^{k}$ divides $2025!$.

Power of $2$ in $2025!$: $1012 + 506 + 253 + 126 + 63 + 31 + 15 + 7 + 3 + 1 = 2017$.\\
Power of $3$ in $2025!$: $675 + 225 + 75 + 25 + 8 + 2 = 1010$.

We need $2k \le 2017$ and $k \le 1010$, so $k \le \min(1008, 1010) = 1008$.
\[
N_{12} = 1008, \quad N_{10} - N_{12} = 505 - 1008 = -503.
\]
\textbf{Answer: (A) $-503$.}

\textit{Why CAT-level:} The non-prime base introduces the constraint $2^{2k}$, not $2^{k}$; the natural error is to assume the limiting factor is $3$.

\vspace{4pt}
\textbf{Solution 13.} \textit{[Area ratio via cevians]}\\
Set $A = (0, 6)$, $B = (0, 0)$, $C = (5, 0)$ (any convenient triangle works; the ratio is invariant). Then
\[
D = (2, 0) \text{ on } BC,\quad E = (5/3, 4) \text{ on } AC.
\]
Line $AD$ parametrised as $(2t,\ 6-6t)$; line $BE$ as $(5s/3,\ 4s)$. Solving:
\[
t = \tfrac{5s}{6},\quad 6(1-t) = 4s \Rightarrow s = \tfrac{2}{3},\ t = \tfrac{5}{9}.
\]
So $P = \big(\tfrac{10}{9},\ \tfrac{24}{9}\big)$ and $AP : PD = 5 : 4$.

Now compute areas. $[ABC] = \tfrac{1}{2}\cdot 5\cdot 6 = 15$.
\[
[APE] = \tfrac{1}{2}\Big| x_A(y_P - y_E) + x_P(y_E - y_A) + x_E(y_A - y_P) \Big| = \tfrac{1}{2}\cdot\tfrac{270}{81} = \tfrac{5}{3}.
\]
Hence $[APE] : [ABC] = \tfrac{5/3}{15} = \dfrac{1}{9}$.
\textbf{Answer: (B) $1 : 9$.}

\textit{Why CAT-level:} Mass points easily give $AP : PD$, but obtaining the area of $\triangle APE$ requires another mass-point or vector step -- the trap is stopping at the $AP : PD$ ratio.

\vspace{4pt}
\textbf{Solution 14.} \textit{[Incircle touch point to midpoint of opposite side]}\\
With $a = BC = 13$, $b = CA = 15$, $c = AB = 14$, semi-perimeter $s = 21$. Tangent length from $A$: $AP = s - a = 8$.

\textbf{Median length.} By Apollonius for median $AM$ from $A$ to mid-$BC$:
\[
AM^{2} = \frac{2b^{2} + 2c^{2} - a^{2}}{4} = \frac{2(225) + 2(196) - 169}{4} = \frac{673}{4}.
\]

\textbf{Angle $\angle BAM$.} Apply the law of cosines in $\triangle ABM$, with $BM = 13/2$:
\[
BM^{2} = AB^{2} + AM^{2} - 2\,AB\cdot AM\cos\angle BAM
\Rightarrow \tfrac{169}{4} = 196 + \tfrac{673}{4} - 14\sqrt{673}\cos\angle BAM,
\]
\[
\Rightarrow 14\sqrt{673}\cos\angle BAM = 322 \Rightarrow \cos\angle BAM = \tfrac{23}{\sqrt{673}}.
\]

\textbf{Apply law of cosines in $\triangle AMP$:} $AP = 8$, $AM = \tfrac{\sqrt{673}}{2}$:
\[
MP^{2} = 64 + \tfrac{673}{4} - 2 \cdot 8 \cdot \tfrac{\sqrt{673}}{2} \cdot \tfrac{23}{\sqrt{673}}
       = 64 + \tfrac{673}{4} - 184 = \tfrac{193}{4}.
\]
So $4\,MP^{2} = 193$.
\textbf{Answer: 193.}

\textit{Why CAT-level:} Coordinate geometry works but produces fractions with denominator $169$. The slick path requires combining median formula + law of cosines, an approach not directly signalled by the problem.

\vspace{4pt}
\textbf{Solution 15.} \textit{[Stacked spheres in a cone]}\\
Place the apex of the cone at the origin and let depth $d$ increase downward along the axis. A sphere of radius $\rho$ centred on the axis at depth $d$ is tangent to the lateral surface iff $\rho = d \sin\alpha$.

For $S_{1}$: $\rho_{1} = 3$, $\sin\alpha = 3/5 \Rightarrow d_{1} = 3 / (3/5) = 5$.

Let $S_{2}$ have radius $\rho_{2}$ centred at depth $d_{2} < d_{1}$. Tangency to lateral surface:
$\rho_{2} = d_{2} \sin\alpha = \tfrac{3}{5}d_{2}$. External tangency to $S_{1}$:
\[
d_{1} - d_{2} = \rho_{1} + \rho_{2} \;\Longrightarrow\; 5 - d_{2} = 3 + \tfrac{3}{5}d_{2}.
\]
\[
\Rightarrow 2 = \tfrac{8}{5}d_{2} \Rightarrow d_{2} = \tfrac{5}{4},\qquad \rho_{2} = \tfrac{3}{5}\cdot\tfrac{5}{4} = \tfrac{3}{4}.
\]
\textbf{Answer: 0.75 cm.}

\textit{Why CAT-level:} The compact relation $\rho = d \sin\alpha$ is not standard CAT vocabulary; deriving it via the cross-sectional triangle and then using the external-tangency condition is what makes this Very Hard.

\hrule

\section*{Topic-Wise Distribution}

\begin{center}
\begin{tabular}{|l|c|c|c|c|c|}
\hline
\textbf{Topic} & \textbf{Total} & \textbf{Medium-Hard} & \textbf{Hard} & \textbf{Very Hard} & \textbf{MCQ / TITA} \\
\hline
Arithmetic         & 5 & 2 & 3 & 0 & 3 / 2 \\
Algebra            & 4 & 0 & 3 & 1 & 3 / 1 \\
Number System      & 3 & 1 & 2 & 0 & 2 / 1 \\
Geometry / Mensuration & 3 & 0 & 1 & 2 & 1 / 2 \\
\hline
\textbf{Total}     & \textbf{15} & \textbf{3} & \textbf{9} & \textbf{3} & \textbf{9 / 6} \\
\hline
\end{tabular}
\end{center}

\hrule

\section*{Quality Verification Summary}

\begin{center}
\begin{tabular}{|c|c|c|c|c|c|c|}
\hline
\textbf{Q.} & \textbf{Answer Verified} & \textbf{MCQ Options Verified} & \textbf{Ambiguity Check} & \textbf{Brute Force Resistance} & \textbf{Method Selection Needed} & \textbf{Status} \\
\hline
1  & Yes & Yes            & Clear wording & Moderate & Yes & Ready \\
2  & Yes & Not Applicable & Clear wording & Yes & Yes & Ready \\
3  & Yes & Yes            & Clear wording & Yes & Yes & Ready \\
4  & Yes & Yes            & Clear wording & Moderate & Yes & Ready \\
5  & Yes & Not Applicable & Clear wording & Yes & Yes & Ready \\
6  & Yes & Yes            & Clear wording & Yes & Yes & Ready \\
7  & Yes & Yes            & Clear wording & Yes & Yes & Ready \\
8  & Yes & Not Applicable & Clear wording & Yes & Yes & Ready \\
9  & Yes & Yes            & Clear wording & Yes & Yes & Ready \\
10 & Yes & Not Applicable & Clear wording & Yes & Yes & Ready \\
11 & Yes & Yes            & Clear wording & Moderate & Yes & Ready \\
12 & Yes & Yes            & Clear wording & Yes & Yes & Ready \\
13 & Yes & Yes            & Clear wording & Yes & Yes & Ready \\
14 & Yes & Not Applicable & Clear wording & Yes & Yes & Ready \\
15 & Yes & Not Applicable & Clear wording & Yes & Yes & Ready \\
\hline
\end{tabular}
\end{center}

\end{document}
